\origpage{29}{}
\textit{Subproof.} Since $r \ge 4$, it follows that $c(\Sigma) \ge 1$. If $c(\Sigma) \ge 2$ then\fix{, since $\Gamma$ is a forest (because $L$ is exhaustive),} there is a region of $\Gamma$ incident with vertices in two different cuffs,\fixnote{fix:6.3-claim3}{For a general drawing this is false: in an annulus, a circuit of $\Gamma$ around the hole puts every region on one side of it. The step needs that $L$ is exhaustive, which makes $\Gamma$ a forest, as stated on page 24.} and so there is a proper line $F$ with ends vertices of $\Gamma$ in different cuffs, and with no internal point in $U(\Gamma)$. But then $F$ does not separate $\Sigma$, contrary to (2). This proves (3).

\hspace*{12pt}(4) \textit{We may assume that for every proper line $F \subseteq \Sigma$ with ends in $V(\Gamma)$ and with no internal point in $U(\Gamma)$, there is a closed disc $\Delta \subseteq \Sigma$ with $F \subseteq bd(\Delta) \subseteq F \cup bd(\Sigma)$.}

\textit{Subproof.} By (2), $F$ separates $\Sigma$. Let $\Sigma_1, \Sigma_2 \subseteq \Sigma$ be surfaces with $\Sigma_1 \cup \Sigma_2 = \Sigma$ and $\Sigma_1 \cap \Sigma_2 = F$. For $i = 1, 2$ let $\Gamma_i = \Gamma \cap \Sigma_i$, and let $G_i = \Gamma_{\fix{i}} \cup \mu(bd(\Sigma) \cap \Sigma_i)$.\fixnote{fix:6.3-claim4-Gi}{Here $G_i$ is built from $\Gamma_i$, the part of $\Gamma$ in $\Sigma_i$, as the next sentence (``using $\Gamma_i$'') needs.} Now for $i = 1, 2$, there is a presentation of $G_i$ in $\Sigma_i$ of depth $\le q$ and defect $\le r + 1$, using $\Gamma_i$; and $L \cap G_i$ is an exhaustive vital $(p + 2q)$-linkage\fix{, except when both ends of $F$ are in one support; in that case, after the vertices of $V(\mu(a)) \cap V(\mu(b))$ on that side are split as in (2), by (2.4), there is such a presentation of the new graph and an exhaustive vital $(p + 4q)$-linkage in it}.\fixnote{fix:6.3-claim4-support}{When both ends $a, b$ of $F$ are in one support $F_j^*$, the side $\Sigma_i$ meets $F_j^*$ in two arcs. The graphs of both arcs contain $V(\mu(a)) \cap V(\mu(b))$, which need not be empty (for instance with $q = 2$). The pieces of a presentation must be disjoint, so the printed claim fails there. After the splitting, each side has at most $\lambda(p + 4q, q, r + 1)$ jumps, so $2q$ becomes $4q$ in the definition of $\lambda$ on page 28 and in the bound below.} If both $\Sigma_1$ and $\Sigma_2$ are not discs, then by (4.2)(i) and (4.2)(ii), either $\hat{\Sigma}_i$ is simpler than $\hat{\Sigma}$, or $\hat{\Sigma}_i$ is homeomorphic to $\hat{\Sigma}$ and $c(\Sigma_i) < c(\Sigma)$ for $i = 1, 2$. But then from (1), there are $\le \lambda(p + \delm{2q}\fix{4q}, q, r + 1)$ $\Sigma_i$-jumps in $\Gamma_i$ for $i = 1, 2$, and hence $\le 2\lambda(p + \delm{2q}\fix{4q}, q, r + 1) \le \lambda$ $\Sigma$-jumps in $\Gamma$, as required. Consequently we may assume that one of $\Sigma_1, \Sigma_2$ is a disc. This proves (4).

From (4), it follows in particular that every $\Sigma$-jump $J$ is homotopic in $\hat{\Sigma}$ to the lines in $bd(\Sigma)$ joining the ends of $J$. Let $R$ be the closure in $\hat{\Sigma}$ of $\hat{\Sigma} \setminus \Sigma$; thus $R \cap \Sigma$ is the unique cuff of $\Sigma$. It follows that every $O$-arc in $R \cup U(\Gamma)$ is null-homotopic in $\hat{\Sigma}$. By theorem (11.10) of [5] there is a closed disc $\Delta \subseteq \hat{\Sigma}$ with $R \cup U(\Gamma) \subseteq \Delta$. By extending $\Delta$ to a sphere and removing $R \setminus bd(R)$ from it, we deduce that $G$ has a presentation in a disc $\Sigma'$, with depth $\le q$ and defect $\le r$. Then the result follows from (1), since $\lambda \ge \lambda(p + 4q, q, r + 2) \ge \lambda(p, q, r)$ and $\Sigma$ is not a disc. $\square$

\sectionhead{7 Presentations and tree-width}

If $\mathcal{T}$ is a tangle in a graph $G$, a subgraph $A$ of $G$ is \textit{small} (with respect to $\mathcal{T}$) if there exists $B \subseteq G$ such that $(A, B) \in \mathcal{T}$. An easy consequence of theorem (2.9) of [6] is the following\del{ (we omit the proof)}.

\textbf{7.1} \textit{If $\mathcal{T}$ is a tangle in $G$, and $A \subseteq G$, then the following are equivalent:}

\hspace*{12pt}(i) \textit{$A$ is small with respect to $\mathcal{T}$;}

\hspace*{12pt}(ii) \textit{if $X$ denotes the set of vertices of $A$ incident with edges of $G$ not $A$, then $|X| < ord(\mathcal{T})$, and $(A, B) \in \mathcal{T}$ where $V(B) = (V(G) \setminus V(A)) \cup X$ and $E(B) = E(G) \setminus E(A)$;}
